\documentclass[aspectratio=169,11pt]{beamer}

% =========================
% Paquetes básicos
% =========================
\usepackage[utf8]{inputenc}
\usepackage[T1]{fontenc}
\usepackage[spanish]{babel}
\usepackage{amsmath,amssymb,amsfonts}
\usepackage{graphicx}
\usepackage{booktabs}
\usepackage{tikz}
\usepackage{xcolor}

% =========================
% Tema visual
% =========================
\usetheme{Madrid}
\usecolortheme{seahorse}
\setbeamertemplate{navigation symbols}{}
\setbeamertemplate{caption}[numbered]

\definecolor{azulEPN}{RGB}{26,76,128}
\definecolor{grisSuave}{RGB}{245,247,250}
\definecolor{naranjaTVq}{RGB}{214,104,49}
\setbeamercolor{title}{fg=azulEPN}
\setbeamercolor{frametitle}{fg=azulEPN}
\setbeamercolor{block title}{bg=azulEPN!85,fg=white}
\setbeamercolor{block body}{bg=grisSuave,fg=black}
\setbeamercolor{alerted text}{fg=naranjaTVq}

% =========================
% Comandos útiles
% =========================
\newcommand{\R}{\mathbb{R}}
\newcommand{\TV}{\operatorname{TV}}
\newcommand{\TVq}{\operatorname{TV}_q}
\newcommand{\grad}{\nabla}
\newcommand{\norm}[1]{\left\lVert #1\right\rVert}
\newcommand{\abs}[1]{\left\lvert #1\right\rvert}

% =========================
% Datos
% =========================
\title{Modelos $TV^q$ no convexos en Restauración de Imágenes}
\subtitle{análisis y un solucionador basado en regularización de región de confianza con convergencia superlineal}
\author{Hintermüller y Wu (2013)}
\institute{MODEMAT}
\date{28 de mayo de 2026}

\begin{document}

% ============================================================
% 1
% ============================================================
\begin{frame}
    \titlepage
\end{frame}

\begin{frame}{Esquema de la presentación}
    \begin{itemize}
    \item Motivación del modelo \(\TV_q\), con \(0<q<1\).
    \item Planteamiento del modelo variacional no convexo.
    \item Regularización tipo Huber del término no diferenciable.
    \item Método Newton primal-dual regularizado.
    \item Resultados teóricos y numéricos.
\end{itemize}
\end{frame}

% ============================================================
% 2
% ============================================================
\begin{frame}{Motivación}
\begin{block}{Pregunta central}
¿Cómo restaurar una imagen degradada preservando bordes y promoviendo la dispersión en el gradiente?
\end{block}
\end{frame}

% ============================================================
% 3
% ============================================================
\begin{frame}{Problema de restauración de imágenes}
\begin{columns}[T]
\column{0.48\textwidth}
\begin{block}{Datos observados}
Se observa una imagen degradada
\[
    z = Ku + \eta,
\]
donde:
\begin{itemize}
    \item \(u\): imagen verdadera,
    \item \(K\): operador de degradación, por ejemplo blur,
    \item \(\eta\): ruido.
\end{itemize}
\end{block}

\column{0.48\textwidth}
\begin{block}{Meta}
Reconstruir una imagen \(u\) que:
\begin{itemize}
    \item sea cercana a los datos \(z\),
    \item sea regular,
    \item conserve bordes importantes.
\end{itemize}
\end{block}
\end{columns}
\end{frame}

% ============================================================
% 4
% ============================================================
\begin{frame}{El modelo clásico: Variación Total}
El enfoque de variación total busca resolver, en forma simplificada,
\[
    \min_u \frac{\lambda}{2}\norm{Ku-z}_2^2 + \alpha \TV(u),
\]
donde
\[
    \TV(u) =\sum_{(i,j)\in\Omega} \abs{(\nabla u)_{ij}}.
\]

\begin{block}{Ventaja principal}
La variación total penaliza el gradiente, pero permite discontinuidades: por eso preserva bordes mejor que regularizaciones cuadráticas.
\end{block}
\end{frame}

% ============================================================
% 5
% ============================================================
\begin{frame}{Dispersión en el gradiente}
En restauración de imágenes, muchas imágenes naturales son aproximadamente constantes por regiones.

\vspace{0.2cm}
\begin{center}
\begin{tikzpicture}[scale=0.9]
    \draw[fill=blue!10,draw=azulEPN,thick] (0,0) rectangle (2,2);
    \draw[fill=orange!15,draw=naranjaTVq,thick] (2,0) rectangle (4,2);
    \draw[very thick,red] (2,0) -- (2,2);
    \node at (1,1) {región 1};
    \node at (3,1) {región 2};
    \node[red] at (2,-0.35) {borde};
\end{tikzpicture}
\end{center}

\begin{block}{Idea}
Si el gradiente es cero en muchas zonas y distinto de cero solo en los bordes, entonces queremos promover sparsity de \(\nabla u\).
\end{block}
\end{frame}

% ============================================================
% 6
% ============================================================
\begin{frame}{De \(\ell_1\) a \(\ell_q\): la idea del paper}
La TV clásica usa una penalización tipo \(\ell_1\) sobre el gradiente.

\[
    \sum_i \abs{v_i}
\]

El paper propone usar una penalización no convexa:
\[
    \sum_i \abs{v_i}^q, \qquad 0<q<1.
\]

\begin{block}{Interpretación}
Para \(0<q<1\), la penalización se acerca más a contar componentes no nulas, es decir, se aproxima mejor a una idea tipo \(\ell_0\).
\end{block}
\end{frame}

% ============================================================
% 7
% ============================================================
\begin{frame}{Contribución del paper}
\begin{enumerate}
    \item Introduce un modelo \(\TV_q\) no convexo para restauración de imágenes.
    \item Prueba existencia de minimizador en el caso discreto.
    \item Regulariza localmente el término no diferenciable mediante una función tipo Huber.
    \item Propone un método Newton primal-dual semismooth.
    \item Usa regularización basada en región de confianza para tratar Hessianos indefinidos.
    \item Demuestra convergencia global y convergencia local superlineal hacia puntos estacionarios del problema regularizado.
\end{enumerate}
\end{frame}

% ============================================================
% 8
% ============================================================
\begin{frame}{Notación básica}
\begin{itemize}
    \item \(\Omega\): dominio discreto de píxeles.
    \item \(u\in\R^{|\Omega|}\): imagen reconstruida.
    \item \(z\in\R^{|\Omega|}\): datos observados.
    \item \(K\in\R^{|\Omega|\times |\Omega|}\): operador de degradación.
    \item \(\nabla u = (\nabla_x u,\nabla_y u)\): gradiente discreto.
    \item \(\lambda_{ij}>0\): peso de fidelidad, posiblemente dependiente del píxel.
\end{itemize}

\begin{block}{Norma local}
Para cada píxel:
\[
    \abs{(\nabla u)_{ij}} = \sqrt{((\nabla_xu)_{ij})^2+((\nabla_yu)_{ij})^2}.
\]
\end{block}
\end{frame}

% ============================================================
% 9
% ============================================================
\begin{frame}{Modelo variacional \(\TV_q\)}
El problema principal en el caso discreto es:
\[
\min_{u\in\R^{|\Omega|}} f(u):=
\sum_{(i,j)\in\Omega}
\left(
\frac{\mu}{2}\abs{(\nabla u)_{ij}}^2
+\frac{\alpha}{q}\abs{(\nabla u)_{ij}}^q
+\frac{\lambda_{ij}}{2}\abs{(Ku-z)_{ij}}^2
\right).
\]

\begin{block}{Parámetros}
\[
    \alpha>0, \qquad 0<q<1, \qquad 0<\mu\ll \alpha.
\]
\end{block}
\end{frame}

% ============================================================
% 10
% ============================================================
\begin{frame}{Planteamiento del modelo}
\begin{columns}[T]
\column{0.32\textwidth}
\begin{block}{Fidelidad}
\[
\frac{\lambda_{ij}}{2}|(Ku-z)_{ij}|^2
\]
Obliga a que \(Ku\) se parezca a los datos observados.
\end{block}

\column{0.32\textwidth}
\begin{block}{TV\(_q\)}
\[
\frac{\alpha}{q}|(\nabla u)_{ij}|^q
\]
Promueve sparsity del gradiente y preserva bordes.
\end{block}

\column{0.32\textwidth}
\begin{block}{Término \(H^1\)}
\[
\frac{\mu}{2}|(\nabla u)_{ij}|^2
\]
Aporta suavidad y estabilidad adicional.
\end{block}
\end{columns}
\end{frame}

% ============================================================
% 11
% ============================================================
\begin{frame}{Existencia de minimizador}
\begin{block}{Hipótesis clave}
\[
    \mathrm{Ker}(\nabla)\cap \mathrm{Ker}(K)=\{0\}.
\]
\end{block}

\begin{block}{Resultado}
Bajo condiciones adecuadas sobre \(\mu,\alpha,q,\lambda\) y \(K\), existe al menos un minimizador global del problema discreto.
\end{block}

\begin{alertblock}{Intuición}
La hipótesis evita direcciones no controladas: no debe existir una imagen constante que también sea anulada por \(K\).
\end{alertblock}
\end{frame}

% ============================================================
% 12
% ============================================================
\begin{frame}{No convexidad y no diferenciabilidad}
Para \(0<q<1\), la función
\[
    s\mapsto |s|^q
\]
es no convexa y no es Lipschitz en cero.

\begin{block}{Conjunto activo e inactivo}
\[
A(u)=\{(i,j)\in\Omega: |(\nabla u)_{ij}|\neq 0\},
\qquad I(u)=\Omega\setminus A(u).
\]
\end{block}

\begin{alertblock}{Consecuencia}
La ecuación de Euler--Lagrange se debe interpretar separando píxeles activos e inactivos.
\end{alertblock}
\end{frame}

% ============================================================
% 13
% ============================================================
\begin{frame}{Ecuaciones de Euler--Lagrange del modelo original}
Una condición de estacionariedad se escribe como:
\[
\begin{cases}
-\mu\Delta u+K^\top\lambda(Ku-z)+\alpha\nabla^\top\left(|\nabla u|^{q-2}\nabla u\right)=0, & (i,j)\in A(u),\\[4pt]
\nabla u=0, & (i,j)\in I(u).
\end{cases}
\]

\begin{block}{Problema computacional}
El término \(|\nabla u|^{q-2}\) puede explotar cuando \(|\nabla u|\to 0\).
\end{block}
\end{frame}

% ============================================================
% 14
% ============================================================
\begin{frame}{Huberización del término \(\TV_q\)}
Para hacer el problema tratable numéricamente, se suaviza localmente:
\[
\varphi_\gamma(s)=
\begin{cases}
\dfrac{1}{q}|s|^q-\left(\dfrac{1}{q}-\dfrac{1}{2}\right)\gamma^q, & |s|>\gamma,\\[6pt]
\dfrac{1}{2}\gamma^{q-2}|s|^2, & |s|\leq \gamma.
\end{cases}
\]

\begin{block}{Interpretación}
Cerca de cero se reemplaza \(|s|^q\) por una expresión cuadrática más regular.
\end{block}
\end{frame}

% ============================================================
% 15
% ============================================================
\begin{frame}{Modelo Huberizado}
El problema regularizado es:
\[
\min_{u\in\R^{|\Omega|}} f_\gamma(u):=
\sum_{(i,j)\in\Omega}
\left(
\frac{\mu}{2}|(\nabla u)_{ij}|^2
+\alpha\varphi_\gamma(|(\nabla u)_{ij}|)
+\frac{\lambda_{ij}}{2}|(Ku-z)_{ij}|^2
\right).
\]

\begin{block}{Ventaja}
El funcional \(f_\gamma\) es continuamente diferenciable.
\end{block}
\end{frame}

% ============================================================
% 16
% ============================================================
\begin{frame}{Gradiente del problema Huberizado}
La condición de primer orden del problema regularizado es:
\[
\nabla f_\gamma(u)=
-\mu\Delta u+K^\top\lambda(Ku-z)
+\alpha\nabla^\top\left(\max(|\nabla u|,\gamma)^{q-2}\nabla u\right)=0.
\]

\begin{alertblock}{Idea clave}
La función \(\max(|\nabla u|,\gamma)\) evita dividir por cantidades cercanas a cero.
\end{alertblock}
\end{frame}

% ============================================================
% 17
% ============================================================
\begin{frame}{Consistencia cuando \(\gamma\to 0^+\)}
\begin{block}{Resultado del paper}
Si \(u_k\) son puntos estacionarios del problema Huberizado y \(\gamma_k\to0^+\), entonces existe una subsucesión que converge a un punto que satisface las ecuaciones de Euler--Lagrange del problema original.
\end{block}

\begin{alertblock}{Mensaje para la exposición}
La Huberización no es solo un truco numérico: está conectada con el problema original cuando el parámetro \(\gamma\) se hace pequeño.
\end{alertblock}
\end{frame}

% ============================================================
% 18
% ============================================================
\begin{frame}{Formulación primal-dual}
Se introduce una variable dual \(p\) asociada al gradiente:
\[
\begin{cases}
-\mu\Delta u + K^\top\lambda(Ku-z)+\alpha\nabla^\top p=0,\\[4pt]
\max(|\nabla u|,\gamma)^{2-q}p=\nabla u.
\end{cases}
\]

\begin{block}{Ventaja}
Esta formulación separa el equilibrio primal de la relación no lineal entre \(p\) y \(\nabla u\).
\end{block}
\end{frame}

% ============================================================
% 19
% ============================================================
\begin{frame}{Newton semismooth primal-dual}
El objetivo es resolver el sistema no lineal primal-dual mediante un método de Newton generalizado.

\begin{block}{Motivación}
Los métodos Newton pueden ser muy rápidos localmente, pero aquí aparecen dos dificultades:
\begin{itemize}
    \item no convexidad del funcional,
    \item posible indefinitud del Hessiano generalizado.
\end{itemize}
\end{block}

\begin{alertblock}{Estrategia del paper}
Combinar Newton semismooth con reweighting y regularización tipo región de confianza.
\end{alertblock}
\end{frame}

% ============================================================
% 20
% ============================================================
\begin{frame}{Reweighting de la ecuación de Euler--Lagrange}
Dado un iterado \(u^k\), se define
\[
    m^k=\max(|\nabla u^k|,\gamma),
\]
y el peso
\[
    w^k=(m^k)^{q+r-2}, \qquad 0\leq r\leq 2-q.
\]

La ecuación reponderada queda:
\[
-\mu\Delta u + K^\top\lambda(Ku-z)
+\alpha\nabla^\top\left(w^k\max(|\nabla u|,\gamma)^{-r}\nabla u\right)=0.
\]
\end{frame}

% ============================================================
% 21
% ============================================================
\begin{frame}{Sistema lineal de Newton}
Después de linealizar y eliminar la variable dual, se obtiene un sistema de la forma:
\[
    \widetilde H_k(r)\,\delta u^{k+1}=-g^k,
\]
donde
\[
    g^k=\nabla f_\gamma(u^k).
\]

\begin{block}{Lectura computacional}
En cada iteración se busca una dirección \(\delta u^{k+1}\) que reduzca el residuo de estacionariedad.
\end{block}
\end{frame}

% ============================================================
% 22
% ============================================================
\begin{frame}{El problema del Hessiano indefinido}
En problemas convexos, el Hessiano suele ser semidefinido positivo.

\vspace{0.2cm}
En este paper, por \(0<q<1\):
\begin{itemize}
    \item el funcional es no convexo,
    \item el Hessiano generalizado puede no ser positivo definido,
    \item una dirección de Newton puede no ser de descenso.
\end{itemize}

\begin{alertblock}{Consecuencia}
Hay que estabilizar el sistema de Newton.
\end{alertblock}
\end{frame}

% ============================================================
% 23
% ============================================================
\begin{frame}{Regularización \(R\)}
El paper interpreta el reweighting como una regularización del Hessiano:
\[
    (H_k+\beta_k R_k)\delta u^{k+1}=-g^k.
\]

\begin{block}{Rol de \(\beta_k\)}
\begin{itemize}
    \item Si \(\beta_k\) es grande, el sistema es más estable.
    \item Si \(\beta_k\to 0\), se recupera el comportamiento de Newton puro.
\end{itemize}
\end{block}

\begin{alertblock}{Idea clave}
Regularizar lejos de la solución, pero dejar que la regularización desaparezca cerca de la solución.
\end{alertblock}
\end{frame}

% ============================================================
% 24
% ============================================================
\begin{frame}{Región de confianza}
Se usa un modelo cuadrático local:
\[
    h_k(d)=f_\gamma(u^k)+(g^k)^\top d+\frac12 d^\top H_k^+d.
\]

La dirección se controla mediante una restricción tipo región de confianza:
\[
    \frac12 d^\top R_{k,\varepsilon}^+d\leq \frac12\sigma_k^2.
\]

\begin{block}{Propósito}
Elegir adaptativamente la cantidad de regularización necesaria para obtener una dirección confiable.
\end{block}
\end{frame}

% ============================================================
% 25
% ============================================================
\begin{frame}{Algoritmo adaptativamente regularizado}
\begin{enumerate}
    \item Inicializar \(u^0,p^0,\beta_0,\sigma_0\).
    \item Construir \(H_k^+\), \(R_{k,\varepsilon}^+\) y \(g^k\).
    \item Resolver
    \[
        (H_k^+ + \beta_k R_{k,\varepsilon}^+)d^k=-g^k.
    \]
    \item Verificar si \(d^k\) es dirección de descenso.
    \item Actualizar \(\beta_k\) y \(\sigma_k\) con criterio de región de confianza.
    \item Hacer búsqueda lineal de Wolfe--Powell.
    \item Actualizar \(u^{k+1}=u^k+a_kd^k\).
\end{enumerate}
\end{frame}

% ============================================================
% 26
% ============================================================
\begin{frame}{Convergencia global}
\begin{block}{Resultado}
Bajo las condiciones del algoritmo y usando búsqueda lineal de Wolfe--Powell,
\[
    \lim_{k\to\infty}\nabla f_\gamma(u^k)=0.
\]
\end{block}

\begin{alertblock}{Interpretación}
El algoritmo no necesariamente encuentra un mínimo global, porque el problema es no convexo, pero sí converge hacia un punto estacionario del problema Huberizado.
\end{alertblock}
\end{frame}

% ============================================================
% 27
% ============================================================
\begin{frame}{Convergencia local superlineal}
\begin{block}{Resultado local}
Si la sucesión converge a un punto estacionario adecuado y los elementos del B-Hessiano son definidos positivos, entonces:
\[
    \norm{u^{k+1}-u^*}=o\left(\norm{u^k-u^*}\right).
\]
\end{block}

\begin{alertblock}{Idea central}
La regularización \(R\) se apaga cerca de la solución:
\[
    \beta_k\to0.
\]
Así se recupera la velocidad rápida de Newton.
\end{alertblock}
\end{frame}

% ============================================================
% 28
% ============================================================
\begin{frame}{Resultados numéricos del paper}
Los autores prueban el método en problemas de:
\begin{itemize}
    \item denoising,
    \item deblurring,
    \item comparación con TV convexa,
    \item sensibilidad al valor inicial.
\end{itemize}

\begin{block}{Parámetros típicos reportados}
En varios experimentos se usa, por ejemplo,
\[
    q=0.75, \qquad \mu=10^{-4}\alpha, \qquad \gamma=0.1.
\]
\end{block}
\end{frame}

% ============================================================
% 29
% ============================================================
\begin{frame}{Plantilla para mostrar resultados visuales}
\begin{columns}[T]
\column{0.25\textwidth}
\centering
\fbox{\parbox[c][2.3cm][c]{0.85\linewidth}{\centering Original}}

\column{0.25\textwidth}
\centering
\fbox{\parbox[c][2.3cm][c]{0.85\linewidth}{\centering Ruidosa}}

\column{0.25\textwidth}
\centering
\fbox{\parbox[c][2.3cm][c]{0.85\linewidth}{\centering TV clásica}}

\column{0.25\textwidth}
\centering
\fbox{\parbox[c][2.3cm][c]{0.85\linewidth}{\centering TV\(_q\)}}
\end{columns}

\vspace{0.4cm}
\begin{alertblock}{Mensaje que se debe enfatizar}
El modelo \(\TV_q\) tiende a preservar mejor los bordes y a promover mayor sparsity del gradiente que la TV convexa.
\end{alertblock}
\end{frame}

% ============================================================
% 30
% ============================================================
\begin{frame}{Conclusiones y bibliografía mínima}
\begin{block}{Conclusiones}
\begin{itemize}
    \item \(\TV_q\), con \(0<q<1\), es una alternativa no convexa a la TV clásica.
    \item La Huberización permite tratar numéricamente la no diferenciabilidad cerca de cero.
    \item El método Newton primal-dual es estabilizado mediante regularización \(R\) y región de confianza.
    \item El esquema combina convergencia global con convergencia local superlineal.
    \item Para tesis sobre TV cuasinormada, este paper es una referencia central.
\end{itemize}
\end{block}

\vspace{0.1cm}
{\footnotesize
\textbf{Referencia:} M. Hintermüller y T. Wu (2013). \emph{Nonconvex TVq-Models in Image Restoration: Analysis and a Trust-Region Regularization-Based Superlinearly Convergent Solver}. SIAM Journal on Imaging Sciences, 6(3), 1385--1415.
}
\end{frame}

\end{document}
